\section{Architecture of AER-MQoS}
\label{sec:arch}

This section makes the \emph{modular boundaries}, \emph{data-path versus control-path responsibilities}, and \emph{Cooja-based validation scope} of AER-MQoS explicit: (i)~implemented mechanisms in the Contiki-NG tree versus forward-looking hooks; (ii)~how multi-class traffic is preserved at route selection and at packet egress; (iii)~no claim of large-scale field trials. Everything labeled ``implemented'' refers to the open-source tree compiled and exercised under Cooja/Contiki-NG unless marked as a design trajectory~\cite{oikonomou2022contik,osterlind2006cooja}.
\textbf{Naming.} This manuscript refers to the protocol exclusively as \emph{AER-MQoS} and documents its C API with the \texttt{aer\_rpl\_plus\_*} symbol prefix and the objective-function source \texttt{rpl-of-aer-plus.c}.

% Notation reference (inserted once in Section~\ref{sec:arch}).
\begin{aertablecol}[t]
  \caption{Principal notation.}
  \label{tab:notation}
  \begingroup\setlength{\tabcolsep}{2.5pt}%
  \aertabwrapcol{%
  \begin{tabularx}{\linewidth}{@{}>{\RaggedRight\arraybackslash}p{1.35cm}X@{}}
    \toprule
    \textbf{Symbol} & \textbf{Meaning} \\
    \midrule
    AER-MQoS / \texttt{AER-MQoS} & Paper term / firmware tag \\
    $\mathrm{DLY}_{\mathrm{n}}$ & ETX-derived delay \emph{proxy} in MCS (not measured latency) \\
    $C0$--$C3$ & Traffic classes ($C3$ most urgent) \\
    $\gamma$, $\alpha$, $\beta$ & Context coeff.; QoS / energy weights ($\alpha{+}\beta{=}1$) \\
    NRE & Normalized residual energy $[0,100]$ \\
    MCS & Multi-criteria rank score \\
    $\mathrm{ETX}_{\mathrm{n}}$ & Normalized ETX in MCS \\
    PRED, LRS & Predicted energy; link-reliability score (from neighbor ETX) \\
    $\bar{J}_c$ & Mean sink inter-arrival gap (class $c$) \\
    \bottomrule
  \end{tabularx}%
  }%
  \endgroup
\end{aertablecol}


Figure~\ref{fig:arch} summarizes the reference architecture (Table~\ref{tab:notation}).

\begin{figure}[tp]
  \centering
  \includegraphics[width=\linewidth]{Fig_1_Architecture_AER_MQoS}
  \caption{\CaptionFigOne}
  \label{fig:arch}

  \noindent Fig.~\ref{fig:arch} shows the QoS plane, trust/energy context plane, Multi-Criteria Score (MCS) computation, and integration with the objective function. Module boundaries match the Cooja firmware tree.
\end{figure}

\subsection{Design principles}
AER-MQoS adheres to three principles for scientific traceability:
\begin{enumerate}
  \item \textbf{Standards-first control plane.} The DODAG lifecycle, rank semantics, and objective-function interface follow RPL and its companion specifications on routing metrics~\cite{rfc6550,rfc6551,rfc9010}. The implementation targets the stock \texttt{rpl-lite} component of Contiki-NG.
  \item \textbf{Explicit cross-layer contracts.} QoS-sensitive behavior is split between (a) rank shaping through the MCS at parent selection time and (b) bounded application-layer scheduling that determines which copy of user data reaches UDP first. This separation decouples MAC queues from RPL ranks. Note that WRR operates above UDP and does not address head-of-line blocking at the MAC/radio layer (see Section~\ref{sec:dio}).
  \item \textbf{Reproducible artifacts.} All quantitative claims are tied to versioned simulation scripts, fixed random seeds, and exported logs.
\end{enumerate}

\subsection{Functional planes on each node}
Each router hosts three cooperating \emph{functional planes} coordinated by a lightweight decision core:
\begingroup\sloppy
\begin{itemize}
  \item \textbf{QoS differentiation plane.} Exposes four traffic classes $C0$--$C3$ aligned with the C enumeration \texttt{aer\_traffic\_class\_t}. Applications register the active class before invoking routing-sensitive APIs so that the MCS and objective function observe the same urgency the application intends~\cite{rfc6551}.
  \item \textbf{Security and trust plane.} Maintains a compact trust score derived from link statistics (e.g., smoothed ETX) and exposes it as an input to the MCS through a normalized ``risk'' term. This signal is an \emph{indicator} for parent ranking and instability detection, not a substitute for cryptographic authentication at the link layer.
  \item \textbf{Energy context plane.} Tracks normalized residual energy (NRE) and a short-horizon prediction of depletion risk so that the MCS can reduce reliance on stressed parents before failures translate into topology repair storms.
\end{itemize}
\endgroup
Harvesting-specific models (e.g., detailed solar traces) remain configurable hooks in the reference code; the Cooja-based study focuses on reproducible energy \emph{proxies} rather than on uncontrolled outdoor campaigns.

\subsection{Path levels and traffic-class mapping}
\label{sec:arch-pathlevels}
To operationalize service differentiation without rewriting MAC schedulers, AER-MQoS maps each traffic class to one of four abstract path levels (levels~1--4, with level~1 being the most protected). The mapping is deterministic in the reference implementation (\texttt{aer\_rpl\_plus\_choose\_path\_level}): $C3\!\rightarrow\!1$, $C2\!\rightarrow\!2$, $C1\!\rightarrow\!3$, $C0\!\rightarrow\!4$. In the evaluated tree, this value is threaded through the application and UDP payload for tracing and class-consistent context updates; it is not consumed inside \texttt{rpl-of-aer-plus.c} for parent exploration, so the levels should be read as an application-visible contract rather than as an additional RPL control-plane knob unless a future revision wires them into the OF.

\begin{figure}[tp]
  \centering
  \includegraphics[width=\linewidth]{Fig_2_DODAG_Path_Levels}
  \caption{\CaptionFigTwo}
  \label{fig:dodag-levels}

  \noindent Fig.~\ref{fig:dodag-levels} shows an illustrative RPL DODAG (not the 25-node Cooja layout). DIO (down), DAO (up), preferred parent path to the sink/LBR. Class-to-level mapping $C3{\rightarrow}1,\ldots,C0{\rightarrow}4$ is in Section~\ref{sec:arch-pathlevels}.
\end{figure}

\subsection{Context fusion and the \texorpdfstring{$(\alpha,\beta,\gamma)$}{alpha/beta/gamma} model}
AER-MQoS couples domain knowledge and traffic urgency through two nonnegative weights $\alpha$ (QoS emphasis) and $\beta$ (energy emphasis) with $\alpha+\beta=1$ in the real domain. A bounded context coefficient $\gamma\in[0,100]$ (stored as $\gamma_{\times100}$ in the implementation) is computed from (i) the configured application domain (health, industry, agriculture, environment) and (ii) the active traffic class, with additional bias when the battery indicator falls below a critical threshold. The implementation then sets $\alpha=\gamma$ and $\beta=1-\gamma$ using the same hundredths-scale ($\times100$) fixed-point convention as $\gamma$, i.e., $\alpha_{\times100}=\gamma_{\times100}$ and $\beta_{\times100}=100-\alpha_{\times100}$; throughout the manuscript, $\gamma$ denotes this conceptual coefficient. A bounded table-driven nudge can adjust $\gamma$ within a narrow interval (Section~\ref{sec:ql}); the evaluation does not treat it as an independently measured performance lever.

\begin{figure}[tp]
  \centering
  \includegraphics[width=0.88\linewidth]{Fig_3_Context_Weights_Alpha_Beta_Gamma}
  \caption{\CaptionFigThree}
  \label{fig:context-weights}

  \noindent Fig.~\ref{fig:context-weights} plots the per-class MCS weights ($\alpha,\beta$) used to balance QoS and energy terms in (3).
\end{figure}

\subsection{Multi-Criteria Score (MCS)}
The MCS aggregates heterogeneous indicators into a scalar suitable for comparison during parent selection. Let $\mathrm{ETX}_{\mathrm{n}}$ and $\mathrm{DLY}_{\mathrm{n}}$ denote normalized ETX and an \emph{ETX-derived delay surrogate} in $[0,100]$ (not a measured queueing or end-to-end delay). Define a QoS partial score
\[
  \mathrm{MCS}_{\mathrm{QoS}} = \frac{w_1\,\mathrm{ETX}_{\mathrm{n}} + w_2\,\mathrm{DLY}_{\mathrm{n}}}{w_1+w_2},
\]
with implementation weights $(w_1,w_2)$ corresponding to \texttt{W1\_ETX} and \texttt{W2\_DELAY}. An energy partial score combines normalized residual energy, predicted energy, and trust-derived risk,
\[
  \mathrm{MCS}_{\mathrm{En}} = \frac{w_3\,\mathrm{NRE} + w_4\,\mathrm{PRED} + w_5\,(100-\mathrm{TRUST})}{w_3+w_4+w_5},
\]
using $(w_3,w_4,w_5)$ associated with \texttt{W3\_NRE}, \texttt{W4\_PRED\_ENERGY}, and \texttt{W5\_TRUST\_RISK}. The final MCS is the convex mixture $\mathrm{MCS} = (\alpha\,\mathrm{MCS}_{\mathrm{QoS}} + \beta\,\mathrm{MCS}_{\mathrm{En}})$ computed in integer arithmetic as in \texttt{aer\_rpl\_plus\_compute\_mcs()}. Each term maps to a named field in the repository; $\mathrm{DLY}_{\mathrm{n}}$ is computed from smoothed ETX only and must not be confused with application-layer end-to-end latency reported in Section~\ref{sec:eval}.

\subsection{Objective function interface}
The dedicated objective function uses a \textbf{locally-assigned experimental Objective Code Point (OCP~8)} used only for simulation purposes; this is \textbf{not} an IANA-registered code point. It lives in the private-use / experimental space of the RPL DAG Metric container~\cite{rfc6550} and translates the MCS into RPL rank increments while inheriting MRHOF-style hysteresis so that rank changes remain damped under variable link quality~\cite{rfc6719}. It is compiled only when \texttt{rpl-of-aer-plus.c} is part of the firmware image (\texttt{RPL\_MQoS}, \texttt{RPL\_AER}, \texttt{AER-MQoS}); the \texttt{RPL\_STANDARD} tag keeps upstream MRHOF without this translation layer. The OCP is suitable for controlled simulation campaigns and must be officially registered for interoperability if deployment outside a \emph{controlled laboratory or simulation} context is ever contemplated~\cite{rfc6550}.

\subsection{Application-layer weighted round-robin queues}
The module \texttt{aer\_qos\_queue} implements four logical queues---one per traffic class---backed by a fixed-size memory pool and Contiki's list primitives. Enqueue operations copy payloads up to a configurable maximum and associate each entry with the active UDP connection and destination set through a small companion API. A weighted round-robin scheduler with service ratios $6{:}3{:}2{:}1$ for $C3{:}C2{:}C1{:}C0$ drains queues on a periodic timer; each successful dequeue re-applies \texttt{aer\_rpl\_plus\_set\_app\_traffic\_class()} before transmission so that the MCS and OF observe the class of the packet that actually seizes the channel. This design mitigates head-of-line blocking between classes while keeping RAM costs explicit, which was singled out as important for constrained hardware targets even when experiments remain Cooja-based.

\subsection{DIO extensions and interoperability}
\label{sec:dio}
\textbf{Scope.} RFC~6550 DIO/DAO messages carry ranks, objective codes, DAG identifiers, and optional \emph{known} TLV-defined metrics~\cite{rfc6550,rfc6551}. In the \emph{evaluated} Contiki-NG artifacts of this repository, \textbf{no custom DIO or DAO TLV authored by this work is serialized on the wire}: trust, residual energy, $\gamma$, traffic-class semantics, and the microscopic Q-learning bias are \emph{local decision variables} (or derived from neighbor link statistics already maintained by \texttt{rpl-lite}), not additional type-length-value fields appended to ICMPv6 RPL control messages. Parent selection therefore consumes information that is either (i) already present in standard neighbor records (ETX via \texttt{rpl\_neighbor\_get\_link\_stats()}) or (ii) computed on the forwarding node before the objective function maps the MCS to a rank addend (experimental OCP~8 whenever \texttt{rpl-of-aer-plus.c} is linked; native MRHOF otherwise). This subsection lists each signal, states explicitly whether it appears inside a DIO payload today, and points to the implementation hook.

Table~\ref{tab:dio-signals} summarizes the mapping. The ``Implementation locus'' column references the Contiki-NG modules in \texttt{code\_source\_AER\_MQoS/}; the repository path is relative to the Contiki-NG \texttt{examples/AER-MQoS/} tree.

\begin{aertablecol}[t]
  \caption{AER-MQoS signals vs.\ DIO/DAO carriage (\texttt{rpl-of-aer-plus.c} builds). Values are local to the node unless noted; no project-specific TLV is serialized in the evaluated firmware. W.=on-wire; St.=local state.}
  \label{tab:dio-signals}
  \begingroup\sloppy
  \setlength{\tabcolsep}{1.8pt}%
  \renewcommand{\arraystretch}{1.03}%
  \aertabwrapcol{%
  \begin{tabularx}{\linewidth}{@{}>{\RaggedRight\arraybackslash}p{0.95cm}>{\RaggedRight\arraybackslash}p{1.05cm}>{\centering\arraybackslash}p{0.5cm}>{\RaggedRight\arraybackslash}X>{\centering\arraybackslash}p{0.35cm}@{}}
    \toprule
    \textbf{Signal} & \textbf{Role} & \textbf{W.} & \textbf{Module(s)} & \textbf{St.} \\
    \midrule
    Classes $C0$--$C3$ & $\gamma$, WRR order & N & \texttt{aer\_mqos}, \texttt{aer\_qos\_queue} & L \\
    $\gamma,\alpha,\beta$ & MCS weights & N & \texttt{aer\_mqos} & L \\
    ETX & MCS$_{\mathrm{QoS}}$ input & Std-RPL & \texttt{rpl-of-aer-plus} & L \\
    ETX-dly surr. & Rank input & N & \texttt{rpl-of-aer-plus} & L \\
    Link rel. (LRS) & MCS$_{\mathrm{En}}$ risk & N & \texttt{aer\_rpl\_trust} & L \\
    NRE, prediction & MCS$_{\mathrm{En}}$ & N & \texttt{aer\_rpl\_energy} & L \\
    Path level 1--4 & Class bias & N & \texttt{aer\_rpl\_plus} & L \\
    $\gamma$ nudge & Table bias & N & \texttt{aer\_rpl\_qlearn} & I \\
    MCS $\rightarrow$ rank & OCP~8 addend & I & \texttt{rpl-of-aer-plus} & L \\
    \bottomrule
  \end{tabularx}%
  }%
  \endgroup
  \vspace{0.2em}
  \raggedright
  \textbf{W.} (on-wire): N=none, Std-RPL=standard neighbor ETX, I=indirect via rank only; \textbf{St.}: L=local node state.
\end{aertablecol}

\paragraph{Code path (rank computation).}
\begingroup\sloppy
The excerpt below shows how neighbor-local metrics feed the MCS \emph{without} custom DIO fields: the OF rebuilds \texttt{aer\_parent\_metrics\_t} per neighbor, then calls \texttt{aer\_rpl\_plus\_compute\_mcs()}.
\endgroup

\begin{lstlisting}[style=csnippet]
aer_rpl_plus_update_context(
    aer_rpl_plus_get_app_traffic_class(),
    aer_energy_get_nre_x100());
met.etx_x100 = (uint16_t)((etx_clamped * 100UL)
    / LINK_STATS_ETX_DIVISOR);
met.delay_ms = (uint16_t)(80UL + (uint32_t)etx / 16UL);
met.trust_x100 = aer_trust_for_nbr(nb);
met.nre_x100 = aer_energy_get_nre_x100();
met.predicted_energy_x100 = aer_energy_get_pred_x100();
uint16_t mcs = aer_rpl_plus_compute_mcs(&met);
return (uint16_t)(64UL + (uint32_t)mcs * 6UL);
\end{lstlisting}

\begingroup\sloppy
\noindent\textit{Source:} \texttt{mcs\_to\_link\_rank\_addend()} in \texttt{rpl-of-aer-plus.c}.
The OF reads ETX via \texttt{rpl\_neighbor\_get\_link\_stats()}, builds \texttt{met}, and sets \texttt{curr\_instance.mc.type} to \texttt{RPL\_DAG\_MC\_NONE}.
The assignment \texttt{met.delay\_ms}${=}80{+}\mathrm{ETX}/16$ is an \textbf{ETX-derived delay surrogate} for MCS weighting only.
Measured end-to-end latency appears in Fig.~\ref{fig:eval-latency} and Table~\ref{tab:empirical-summary} (TX/RX logs).
\endgroup

\paragraph{Design trajectory (future TLV roadmap).}
The architecture still \emph{reserves} structured fields (trust, energy, anomaly flags) for a potential future DIO option that would advertise capabilities or aggregated health indicators~\cite{rfc6550}. Such TLVs would need IANA-style option codes (or experimental namespaces), on-wire backward compatibility, and a second simulation campaign before any claim of interoperability beyond Cooja. Until that code ships in the same branch as the measurements, the manuscript must treat TLV extensions strictly as a \textbf{design trajectory}, not as part of the evaluated binary.

\subsection{Interoperability summary}
Unknown DIO options must remain ignorable by strict RFC~6550 parsers; any future TLV work will therefore piggyback on optional, ignorable encodings and will be validated only after it is merged and re-simulated. Section~\ref{sec:eval} contrasts how competing protocol families influence observable metrics under identical Cooja conditions despite the absence of custom DIO payloads in the present artifact.

\subsection{Summary}
AER-MQoS combines (i) a transparent MCS with explicit weights, (ii) MRHOF-like hysteresis at the OF, (iii) deterministic class-to-level mapping, and (iv) application-layer WRR queues, all coordinated under a Cooja-based evaluation methodology. Section~\ref{sec:dio} documents control-plane interoperability: present binaries compute composite ranks without custom DIO TLVs. Table~\ref{tab:proto-metric} and Section~\ref{sec:validation-plan} map observables to protocol behaviors and planned ablations. The next sections instantiate trust, energy, and the Q-learning \emph{specification} referenced here (learning is not experimentally validated as a performance lever in the committed Cooja campaigns of Section~\ref{sec:eval}).
